\section*{第 \ref{ch:g12:meiosis-genetic-shuffling} 章——减数分裂与遗传重组}

\begin{solution}{exo:g12:meiosis-genetic-shuffling:1}
\omterm{def:g12:meiosis-genetic-shuffling:meiosis}{二倍体}：每种\omterm{def:g11:cell-cycle-mitosis:chromosome}{染色体}有两条（$2n$）；\omterm{def:g12:meiosis-genetic-shuffling:meiosis}{单倍体}：只有一条（$n$）。
\omterm{def:g12:meiosis-genetic-shuffling:meiosis}{减数分裂}使\omterm{def:g11:cell-cycle-mitosis:chromosome}{染色体}数由 $2n$ 减半为 $n$，并使
\omterm{def:g10:universal-dna:information}{DNA} 含量（复制后为 $2q$）在每个配子中降至 $q/2$。
\end{solution}

\begin{solution}{exo:g12:meiosis-genetic-shuffling:2}
减数第一次分裂使每对\omterm{def:g10:common-ancestry:homology}{同源}\omterm{def:g11:cell-cycle-mitosis:chromosome}{染色体}彼此分离；减数
第二次分裂使每条\omterm{def:g11:cell-cycle-mitosis:chromosome}{染色体}的两条\omterm{def:g11:cell-cycle-mitosis:chromosome}{染色单体}分开。前者
具有减数性质。
\end{solution}

\begin{solution}{exo:g12:meiosis-genetic-shuffling:3}
减数第一次分裂后：4 条\omterm{def:g11:cell-cycle-mitosis:chromosome}{染色体}（每条含两条\omterm{def:g11:cell-cycle-mitosis:chromosome}{染色单体}），\omterm{def:g10:universal-dna:information}{DNA} 为 $q$。
减数第二次分裂后：4 条\omterm{def:g11:cell-cycle-mitosis:chromosome}{染色体}（每条含一条\omterm{def:g11:cell-cycle-mitosis:chromosome}{染色单体}），\omterm{def:g10:universal-dna:information}{DNA} 为 $q/2$。
\end{solution}

\begin{solution}{exo:g12:meiosis-genetic-shuffling:4}
$2^4 = 16$.
\end{solution}

\begin{solution}{exo:g12:meiosis-genetic-shuffling:5}
与仅携带\omterm{def:g11:genes-and-disease:genetic}{隐性}\omterm{def:g10:universal-dna:gene}{等位基因}的个体杂交，使每个后代
都直接反映被测亲本的配子。重组体的
比例可说明两个\omterm{def:g10:universal-dna:gene}{基因}是否位于同一条\omterm{def:g11:cell-cycle-mitosis:chromosome}{染色体}上，
若是，则还可说明二者相距多远。
\end{solution}

\begin{solution}{exo:g12:meiosis-genetic-shuffling:6}
减数第一次分裂：$2q$；两次分裂之间：$q$；配子：$q/2$。若无
中间的 S 期，第二次分裂会使 \omterm{def:g10:universal-dna:information}{DNA} 再减半，这正
使配子所含 \omterm{def:g10:universal-dna:information}{DNA} 达到\omterm{def:g12:meiosis-genetic-shuffling:meiosis}{二倍体}\omterm{def:g10:cells-common-unit:cell}{细胞}的一半。
\end{solution}

\begin{solution}{exo:g12:meiosis-genetic-shuffling:7}
四类数目相等：两对\omterm{def:g10:universal-dna:gene}{基因}不连锁，位于不同\omterm{def:g11:cell-cycle-mitosis:chromosome}{染色体}上。配子
$AB$、$Ab$、$aB$、$ab$，各占四分之一。
\end{solution}

\begin{solution}{exo:g12:meiosis-genetic-shuffling:8}
连锁：两类数目较多（$AB$、$ab$，亲本型）和两类较少（$Ab$、
$aB$，重组型）。重组体为 $85/1000 = 8.5\%$。
\end{solution}

\begin{solution}{exo:g12:meiosis-genetic-shuffling:9}
只有 $AB$ 和 $ab$：发生在两个\omterm{def:g10:universal-dna:gene}{基因}之外的交换互换了
都不含这两个\omterm{def:g10:universal-dna:gene}{基因}的片段，因此 $A/a$ 与 $B/b$ 的组合不变。
两\omterm{def:g10:universal-dna:gene}{基因}之间发生重组，必须在二者之间发生交换。
\end{solution}

\begin{solution}{exo:g12:meiosis-genetic-shuffling:10}
约为千分之 0.8、2.9 和 16，25 与 42 之间相差 20 倍。
\end{solution}

\begin{solution}{exo:g12:meiosis-genetic-shuffling:11}
减数第一次分裂时，两条\omterm{def:g10:common-ancestry:homology}{同源}\omterm{def:g11:cell-cycle-mitosis:chromosome}{染色体}进入同一个\omterm{def:g10:cells-common-unit:cell}{细胞}：产生两个含两份拷贝
的配子和两个不含该\omterm{def:g11:cell-cycle-mitosis:chromosome}{染色体}的配子。减数第二次分裂时，同一条\omterm{def:g11:cell-cycle-mitosis:chromosome}{染色体}的两条\omterm{def:g11:cell-cycle-mitosis:chromosome}{染色单体}
进入同一个\omterm{def:g10:cells-common-unit:cell}{细胞}：产生一个含两份拷贝的配子、一个不含的配子，
另外两个为正常配子。
\end{solution}

\begin{solution}{exo:g12:meiosis-genetic-shuffling:12}
重组率与距离成正比：2\% 表示两\omterm{def:g10:universal-dna:gene}{基因}
非常接近，交换很少落在二者之间；48\% 表示相距很远，
几乎每次\omterm{def:g12:meiosis-genetic-shuffling:meiosis}{减数分裂}中二者之间都发生交换。由于
一次交换可产生 50\% 的重组配子，相距较远的连锁\omterm{def:g10:universal-dna:gene}{基因}
重组的自由程度与独立遗传的\omterm{def:g10:universal-dna:gene}{基因}相同。
\end{solution}

\begin{solution}{exo:g12:meiosis-genetic-shuffling:13}
$AB$、$Ab$、$aB$、$ab$，各占四分之一，这是两对\omterm{def:g10:universal-dna:gene}{基因}
独立分配的结果。\omterm{prop:g12:meiosis-genetic-shuffling:crossingover}{交叉互换}对这些\omterm{def:g10:universal-dna:gene}{基因}无影响：它
只重组同一条\omterm{def:g11:cell-cycle-mitosis:chromosome}{染色体}上的\omterm{def:g10:universal-dna:gene}{等位基因}，而这两对\omterm{def:g10:universal-dna:gene}{基因}位于
不同的\omterm{def:g11:cell-cycle-mitosis:chromosome}{染色体}上。
\end{solution}

\begin{solution}{exo:g12:meiosis-genetic-shuffling:14}
每位亲本在每个\omterm{def:g10:universal-dna:gene}{基因}位点上随机将两个\omterm{def:g10:universal-dna:gene}{等位基因}中的一个
传递给子女；对某一\omterm{def:g10:universal-dna:gene}{基因}而言，两个同胞从某一亲本
获得相同\omterm{def:g10:universal-dna:gene}{等位基因}的概率为 $1/2$。在多个\omterm{def:g10:universal-dna:gene}{基因}上，共享比例
平均为 $1/2$，但对任意一对同胞而言，该比例围绕此值波动。
\end{solution}

\begin{solution}{exo:g12:meiosis-genetic-shuffling:15}
母本植株的克隆体，与母本及彼此在遗传上完全相同。
该植物获得了快速、可靠的增殖：可使在当前环境中有效的\omterm{def:g11:enzymes-and-phenotype:phenotype}{基因型}
迅速保持下来；但它失去了多样性，
而多样性可使部分后代在环境变化或出现新寄生虫时存活。
\end{solution}

\begin{solution}{pb:g12:meiosis-genetic-shuffling:1}
\textbf{1.} 雌性为 $b^+b$、$vg^+vg$，其中 $b^+vg^+$ 位于一条\omterm{def:g11:cell-cycle-mitosis:chromosome}{染色体}上，
另一条为 $b\,vg$（来自其两个纯合亲本）。雄性为 $bb$，
$vg\,vg$。

\textbf{2.} 仅产生 $b\,vg$ 配子：雄性只提供\omterm{def:g11:genes-and-disease:genetic}{隐性}
\omterm{def:g10:universal-dna:gene}{等位基因}，因此每个后代的\omterm{def:g11:enzymes-and-phenotype:phenotype}{表型}都反映雌性的
配子。

\textbf{3.} $b^+vg^+$（灰身正常翅）、$b\,vg$（黑身残翅），
$b^+vg$（灰身残翅）、$b\,vg^+$（黑身正常翅）。

\textbf{4.} 灰身正常翅 42\%，黑身残翅 41\%，灰身残翅 9\%，
黑身正常翅 8\%。

\textbf{5.} 若位于同一条\omterm{def:g11:cell-cycle-mitosis:chromosome}{染色体}上：两类数目较多、两类较少，
而不是 $1:1:1:1$。

\textbf{6.} 灰身正常翅和黑身残翅：由她的两条\omterm{def:g10:common-ancestry:homology}{同源}\omterm{def:g11:cell-cycle-mitosis:chromosome}{染色体}携带的组合，
完整地遗传自其纯合亲本。

\textbf{7.} 灰身残翅和黑身正常翅，
源于前期 I 中两个\omterm{def:g10:universal-dna:gene}{基因}之间的\omterm{prop:g12:meiosis-genetic-shuffling:crossingover}{交叉互换}。

\textbf{8.} $(206 + 185)/2300 \approx 17\%$。两\omterm{def:g10:universal-dna:gene}{基因}之间的\omterm{prop:g12:meiosis-genetic-shuffling:crossingover}{交叉互换}
使四条\omterm{def:g11:cell-cycle-mitosis:chromosome}{染色单体}中的两条成为重组型，因此 17\%
的重组配子意味着约 34\% 的\omterm{def:g12:meiosis-genetic-shuffling:meiosis}{减数分裂}发生了交换。

\textbf{9.} 单次交换产生两条重组型\omterm{def:g11:cell-cycle-mitosis:chromosome}{染色单体}，
两者各为一种类型；两条\omterm{def:g10:common-ancestry:homology}{同源}\omterm{def:g11:cell-cycle-mitosis:chromosome}{染色体}均等地进入配子，
因此两类亲本型配子的数目也相等。

\textbf{10.} 亲本类型现在为灰身残翅和黑身正常翅，
各约 41.5\%；重组类型为灰身正常翅和黑身残翅，各约
8.5\%。

\textbf{11.} 是：两者的重组率都远低于 50\%。

\textbf{12.} $b$——$cn$——$vg$：9 与 8 相加为 17。

\textbf{13.} 区间内发生交换的概率与其长度成正比，
因此相邻区间的频率可以相加：
重组频率衡量的是沿\omterm{def:g11:cell-cycle-mitosis:chromosome}{染色体}的距离。

\textbf{14.} 一次交换在四条\omterm{def:g11:cell-cycle-mitosis:chromosome}{染色单体}中产生两条重组型\omterm{def:g11:cell-cycle-mitosis:chromosome}{染色单体}，
即至多 50\%；额外的交换只会重新组合，但绝不会使
重组比例超过一半。

\textbf{15.} 它位于另一条\omterm{def:g11:cell-cycle-mitosis:chromosome}{染色体}上，或位于同一条\omterm{def:g11:cell-cycle-mitosis:chromosome}{染色体}上相距很远的位置；
这两种情况无法由杂交实验区分。

\textbf{16.} $2^4 = 16$ 种配子；$16 \times 16 = 256$ 种组合。

\textbf{17.} 每次交换都会产生此前从未存在过的\omterm{def:g11:cell-cycle-mitosis:chromosome}{染色体}，
其交换位置在不同的\omterm{def:g12:meiosis-genetic-shuffling:meiosis}{减数分裂}中各不相同；
因此不同\omterm{def:g11:cell-cycle-mitosis:chromosome}{染色体}及配子的数目
没有实际上限。

\textbf{18.} $2^{23} \times 2^{23} \approx \num{7e13}$：组合数
为迄今为止所有人类总数的七百倍，
而这尚未计入\omterm{prop:g12:meiosis-genetic-shuffling:crossingover}{交叉互换}。

\textbf{19.} 同胞的\omterm{def:g10:universal-dna:gene}{等位基因}取自相同的四组
（每位亲本两组），平均共享其中一半；陌生人则取自
不同的组。

\textbf{20.} 17\%，即雌性配子中
在 $b$ 与 $vg$ 之间发生重组的比例，可度量这两个\omterm{def:g10:universal-dna:gene}{基因}之间的距离；
各对\omterm{def:g11:cell-cycle-mitosis:chromosome}{染色体}的自由组合以及每对\omterm{def:g10:common-ancestry:homology}{同源}\omterm{def:g11:cell-cycle-mitosis:chromosome}{染色体}内部的\omterm{prop:g12:meiosis-genetic-shuffling:crossingover}{交叉互换}
使每个配子各不相同。
\end{solution}
